\documentclass[10pt,a4paper]{article}
\usepackage[T1]{fontenc}
\usepackage[utf8]{inputenc}
\usepackage{lmodern}
\usepackage[margin=1.7cm,top=1.3cm,bottom=1.2cm]{geometry}
\usepackage{amsmath}
\usepackage{booktabs}
\usepackage{microtype}
\pagestyle{empty}
\setlength{\parskip}{0.24em}
\setlength{\parindent}{0pt}

\begin{document}

{\large\bfseries Reply to S. Danilov's comments on the mEVP communication study}\\[0.2em]
{\small N. Koldunov, 6 August 2026}

\vspace{0.45em}\hrule\vspace{0.6em}

We built the variant you describe, and it changes the study's conclusion. Written as you write it
--- the ring computation folded into the loops that already exist rather than into kernels of its
own --- the wide halo removes \textbf{58\% of the sea-ice cost} on CORE2 at its operating point,
against 11\% for our version of the same scheme. The two runs compute the same ring and put
identical message and byte counts on the wire; they differ by 360 kernel launches per time step.
Our earlier conclusion, that exactness costs most of the saving, was measuring our implementation.

\begin{center}\small
\begin{tabular}{@{}lrrr@{}}
\toprule
CORE2, one node, four GPUs & model step & \multicolumn{2}{c}{sea-ice cost} \\
\midrule
every sub-cycle (reference)                     & 0.0660\,s          & 15.8\,ms          & --- \\
wide halo, standard form, our kernels           & $-2.4\%$           & 14.1\,ms          & $-10.8\%$ \\
wide halo, divergence form, our kernels         & $+1.1\%$           & 16.3\,ms          & $+3.2\%$ \\
\textbf{wide halo, standard form, your loops}   & $\mathbf{-13.6\%}$ & \textbf{6.7\,ms}  & $\mathbf{-57.6\%}$ \\
\textbf{wide halo, divergence form, your loops} & $\mathbf{-13.6\%}$ & \textbf{6.6\,ms}  & $\mathbf{-58.2\%}$ \\
\bottomrule
\end{tabular}
\end{center}

58\% is the upper end of the range we measured for the delayed exchange, the scheme we rejected
because it puts the domain decomposition into the ice field. If this holds on the other three
meshes, the trade we reported between speed and a clean ice field does not exist; it was a
property of how we had written the ring. Giving up the owner's summation order costs what you
predicted: after a time step of 120 sub-cycles the ice velocity differs from the exact wide halo
by 8 to 15 units in the last place.

\textbf{Point 1, the auxiliary fields.} We instrumented the per-step exchange to report, field by
field, the difference between the owner's arriving value and the value the receiving process
already holds. Nine of the eleven arrive \emph{exactly} equal --- ice velocity, concentration,
thickness, snow, ocean velocity, wind stress --- because FESOM's own halo exchange has already put
them there, as you said. The two right-hand sides do not: they are assembled and area-scaled over
owned nodes only, so a halo node has no correct value at all. But the redundancy stops at the
first ring, and rings $2\ldots K$ have no other source, so what could be removed is 37\% of that
exchange at $K$=2 and 17\% at $K$=4. Since the exchange is 5.4\% of the data the scheme moves and
one message per neighbour per step, removing it would gain under one percent of the traffic and no
messages. We have not implemented the trimmed version, and that is why.

\textbf{Point 1, the kernel structure.} This is the part we could not bound before, and it was
much the larger. Our wide halo adds 362 kernel launches per model step, each on a slower launch
path than the owned kernel it duplicates: the ring kernels carry enough array descriptors to cross
a threshold in Kokkos above which arguments are staged through constant memory. The gap between
consecutive launches is 14.3\,$\mu$s for those against 4.4\,$\mu$s for the owned kernels. Per step
the ring kernels cost 5.9\,ms of inter-kernel idle against 2.5\,ms of arithmetic.

\textbf{The two forms become indistinguishable.} In our implementation the standard and divergence
forms differed by 3.5 percentage points, and the previous note decomposed that difference into a
message term and a byte term. With the ring folded into the existing loops the difference is zero
to four decimals, while the divergence form still ships 40\% fewer numbers. The gap we decomposed
was our kernel structure. On CPU at 512 cores the picture differs and favours your form: there the
ring assembly is arithmetic-bound, and folding the loops removes a recomputation the divergence
form was doing about six times over, worth 3 percentage points, while in the standard form the
same change costs 5.

\textbf{Point 2, extending the computation over the ring.} You were right, and the correction runs
deeper than we allowed. Recomputing the ring does remove the imprint --- averaging $|\Delta$ ice
thickness$|$ on the sub-domain boundaries against the interior over 60 days on fArc gives 0.87,
against 0.85 for two identical runs and 2.52 for the delayed exchange. What we wrote was that the
same change removes most of the saving and that we had found no intermediate that was both cheap
and clean. That was a statement about our kernels, not about the scheme.

\textbf{Open.} The table is CORE2 at its operating point; fArc, DARS and NG5 are queued. The
60-day boundary test is being repeated for this variant and is the gate we hold it to, since it is
no longer bit-identical and cannot be checked the way the exact version was.

\vspace{0.35em}
{\footnotesize All numbers are from the runs described in \texttt{danilov\_mevp\_report.pdf}; the
output and job scripts can be shared.}

\end{document}
